\documentclass[12pt,a4paper]{article}
\usepackage[utf8]{inputenc}
\usepackage[lithuanian]{babel}
\usepackage[T1]{fontenc}
\usepackage{lmodern}
\usepackage{amsmath, amssymb, amsthm}
\usepackage{geometry}
\geometry{left=2cm, right=2cm, top=2.5cm, bottom=2.5cm}
\usepackage{enumitem}
\usepackage{tasks}

\begin{document}

\begin{center}
    {\Large \textbf{Pasiruošimas kontroliniam darbui 2}}\\[0.5cm]
\end{center}
\vspace{0.5cm}

\section*{Užduotys}

\begin{enumerate}
    \item Kiekvienai iš žemiau esančių reiškinių nustatykite jo tipą (racionalusis, iracionalusis, sveikas, trupmeninis).
    \begin{tasks}[label-width=1.2em](3)
        \task $\frac{x^2 - 4}{x+2}$
        \task $\sqrt{2}x + 5$
        \task $\sqrt{x-3} + 2x$
        \task $\frac{5}{x}$
        \task $\frac{2x}{3}+ \sqrt{2}$
        \task $4x ^{2} +3x ^{-1} $
    \end{tasks}

    \item Apskaičiuokite reiškinio reikšmę:
    $$\frac{2^{2026} - 2^{2025}}{2^{2025}}$$

    \item Remdamiesi laipsnio su natūraliuoju rodikliu apibrėžimu, įrodykite savybę: 
    $$x^m : x^n = x^{m-n} \quad \text{(čia $m, n \in \mathbb{N}$)}.$$
    \item Suprastinkite reiškinį ir atsakymą užrašykite taip, kad nebūtų neigiamų laipsnių rodiklių:
    $$ \left(\frac{2x^{-2}y^3}{3x^4y^{-1}}\right)^{-2}$$
    \item Duotas vienanaris: $-4^2 x^3 y^2 z$. Nustatykite šio vienanario koeficientą ir jo laipsnį.

    \item Atlikite nurodytus veiksmus ir nustatykite gauto vienanario koeficientą bei laipsnį. Jei gautas reiškinys nėra vienanaris paaiškinkite kodėl.
    $$(-3a^2b^3)^3 \cdot \left(\frac{1}{9}ab^{-2}\right)^2$$

    \item Suprastinkite reiškinį:
    $$ (2x - 3)(x^2 + 2x - 1) - 2(x-3) ^{3}  $$

    \item Nenaudodami kvailintuvo apskaičiuokite: $$ \frac{102^3 + 48^3}{(5 \sqrt{6} ) ^{2}  } +24 \cdot 204- \sqrt[3]{48 ^{6} } $$  

    \item Įrodykite kubų skirtumo formulę:
    $$ a^3 - b^3 = (a - b)(a^2 + ab + b^2) $$

    \item Suprastinkite:
    $$ (2x+y)^3 - (2x-y)^3 - 2y^3 $$

    \item Duota kvadratinį trinarį užrašę pavidalu $a(x+ b)^2 + c$ apskaičiuokite reiškinio $16abc+396$ reikšmę:
    $$ 2x^2 - 5x + 13 $$


    \item Raskite didžiausią ir mažiausią reiškinio $-3x^2 + 4x + 1$ reikšmę. Su kokia $x$ reikšme šios reikšmės įgyjamos?

    \item Suprastinkite trupmeną:
    $$\frac{a^2 - b^2}{a^2 + 2ab + b^2} : \frac{a-b}{a+b}$$

    \item Raskite $x$ ir $y$ reikšmes, su kuriomis reiškinys $x^2 + y^2 - 4x + 6y + 13$ įgyja savo \textit{mažiausią} reikšmę.

    \item Įrodykite formulę:
    $$a^4+4 b^4=\left(a^2+2 a b+2 b^2\right)\left(a^2-2 a b+2 b^2\right)$$

    \item Išskaidykite daugianarį dauginamaisiais:
    $$ x^4 + 4 $$
\end{enumerate}

\newpage
\section*{Atsakymai}

\begin{itemize}[label=]
    \item \textbf{1.} a) Racionalusis trupmeninis; b) Racionalusis sveikas; c) Iracionalusis; d) Racionalusis trupmeninis; e) Racionalusis sveikas; f) Racionalusis trupmeninis.
    \item \textbf{2.} $1$
    \item \textbf{4.} $\displaystyle \frac{9x^{12}}{4y^8}$
    \item \textbf{5.} Koeficientas: $-16$, laipsnis: $6$
    \item \textbf{6.} Koeficientas: $-\frac{1}{3}$, laipsnis: $13$
    \item \textbf{7.} $19x^2 - 62x + 57$
    \item \textbf{8.} $15300$
    \item \textbf{10.} $24x^2y$
    \item \textbf{11.} $1$
    \item \textbf{12.} Didžiausia reikšmė yra $2\frac{1}{3}$ (įgyjama, kai $x = \frac{2}{3}$). Mažiausios reikšmės nėra.
    \item \textbf{13.} $1$
    \item \textbf{14.} $x = 2$, $y = -3$
    \item \textbf{16.} $(x^2 - 2x + 2)(x^2 + 2x + 2)$
\end{itemize}

\end{document}
